\section{Verdict on the Hypothesis}
\label{sec_conc_verdict}

The hypothesis stated in Section~\ref{sec_intro_hypothesis} consists of
three separate claims.
Each clause is evaluated against the experimental evidence in turn,
before the overall verdict on the hypothesis is given at the end of the
section.

\subsection{Clause~1: pose grouping is an embedding-quality problem}
\label{sec_conc_verdict_embedding}

\textbf{Supported:}
The first clause proposes that pose-grouping performance is constrained
primarily by the quality of the per-keypoint embedding rather than by the
choice of grouping algorithm applied afterwards.
The experiments provide consistent evidence for this claim.
Three learned grouping heads (DEC,
slot attention,
and graph partitioning) were evaluated on the standard GAT backbone
under matched training conditions (Section~\ref{sec_res_learned_heads}),
followed by ten SA-DMoN variants that explored a separate family of
graph-structural grouping objectives (Section~\ref{sec_res_sadmon}).
None produced a configuration that outperformed plain $k$-means once
transfer to COCO was considered.

The graph-partitioning result is particularly illustrative.
It reaches $0.9714$ synthetic-validation PGA,
making it the strongest of the learned heads on the synthetic
distribution,
but its own grouping read-out falls to $0.7514$ on COCO.
By comparison, the no-depth standard GAT with $k$-means reaches
$0.8841$ on the same COCO evaluation setting.
The SA-DMoN experiments show the same broader pattern from a different
direction -
changing the grouping objective and its graph construction does not
recover the performance lost when the underlying embedding is not
sufficiently transferable.

The embedding-side experiments produce the opposite result.
The three skeleton-aware PG-GAT modifications of
Section~\ref{sec_meth_pggat} improve the representation at matched
architecture size,
the architecture sweep of Section~\ref{sec_res_pggat_sweep} raises
synthetic-validation PGA from $0.8896$ to $0.9634$,
and subsequent COCO fine-tuning raises COCO val PGA from $0.9010$
to $0.9715$.
These gains are substantially larger than those obtained by changing
the grouping head while leaving the basic embedding formulation
unchanged.

Taken together with the analysis of Section~\ref{sec_disc_embedding},
the evidence supports the first clause of the hypothesis.
For the pose-grouping setting studied here,
the main performance lever is the representation in which the detected
keypoints are embedded.
Once that representation separates people effectively,
a simple $k$-means read-out is sufficient -
increasing the complexity of the grouping stage does not provide a
corresponding improvement.

\subsection{Clause~2: skeleton-aware GAT with two-stage training}
\label{sec_conc_verdict_architecture}

\textbf{Supported, with qualification:}
The second clause proposes that the modular grouping approach can achieve
quality competitive with a jointly trained end-to-end alternative.
The architecture and fine-tune sweeps produce the
\texttt{meng-headline} checkpoint (W\&B run \texttt{b5noyg7t}),
which reaches $0.9715$ COCO validation PGA when evaluated on
ground-truth keypoints and $0.9591$ end-to-end PGA when applied to
HigherHRNet detections.
On the same matched detections,
HigherHRNet's native associative-embedding grouping reaches $0.9847$,
leaving a difference of $0.0256$ PGA.

The modular method therefore does not match or exceed the jointly
trained reference on this evaluation.
The result instead establishes a relatively small measured cost for
separating the grouping stage from the detector.
This distinction matters because the two methods provide different
capabilities.
HigherHRNet's AE representation is trained jointly with its detector and
achieves the higher grouping accuracy,
whereas PG-GAT can be applied as a separate grouping module to
detections produced by different architectures.
The detector-swap experiments strengthen this interpretation by showing
a similar autonomous gap of approximately $0.02$ PGA across
HigherHRNet,
YOLO11x-pose,
and RTMO-L.

The hypothesis described the target as "quality competitive with an
end-to-end joint-trained alternative" without defining a numerical
equivalence margin in advance.
The verdict must therefore remain qualitative rather than being
presented as a statistical non-inferiority result.
Within that limitation, a difference of $0.0256$ against a reference
operating at $0.9847$ PGA is considered sufficiently small to support
the intended claim of competitive grouping quality,
particularly in settings where detector independence is itself a
requirement.

The qualification is that this conclusion depends on what the deployment
setting values.
If the objective is solely to maximise PGA with a fixed HigherHRNet
detector, the native AE grouping remains the better choice by $0.0256$.
If grouping must remain independent of the detector,
the experiments show that PG-GAT provides that capability while
retaining $0.9591$ PGA on the same detections.
The second clause is therefore supported as a claim of competitive
rather than superior performance,
with the measured accuracy difference forming part of the modularity
trade-off rather than being hidden by the verdict.

\subsection{Clause~3: autonomous deployment and detector swap}
\label{sec_conc_verdict_autonomy}

\textbf{Supported, with the autonomy result stronger than originally expected:}
The third clause concerns whether the proposed grouping pipeline can
operate without ground-truth knowledge of the number of people and can
be transferred between keypoint detectors without retraining.
Both parts are supported by the experiments.

The autonomous evaluation of Section~\ref{sec_res_e2e} replaces the
ground-truth person count with the prediction of the learned K-head.
The head achieves $76.1\%$ exact-count accuracy and $91.9\%$ off-by-one
accuracy on COCO validation.
Despite this imperfect agreement with the annotated person count,
the resulting pipeline reaches $0.9644$ end-to-end PGA,
compared with $0.9591$ when $k$-means is supplied with the ground-truth
$K$.
Predicted $K$ therefore improves the measured PGA by $0.0053$ rather
than introducing the degradation that might normally be expected when an
oracle quantity is replaced by an estimate.

As discussed in Section~\ref{sec_disc_modularity_autonomy_autonomy},
this result follows from the distinction between the annotated number of
people and the number represented by an incomplete detection set.
The K-head predicts from the detected embeddings and keypoint count,
allowing it to adapt to the evidence that actually reaches the grouping
stage.
Ground-truth $K$ describes the complete annotated image and can
therefore overestimate the appropriate number of clusters when one or
more people are poorly represented by the detector.
The detector-swap experiments and ground-truth-keypoint control support
this interpretation -
the predicted-$K$ advantage grows as detection recall decreases and
reverses when complete ground-truth keypoints are supplied.
The autonomy result is therefore stronger than the original requirement
that predicted $K$ should merely preserve acceptable grouping quality.

The detector-independence part of the clause is supported both by the
architecture and by direct evaluation.
PG-GAT operates on detected keypoints rather than detector-specific
intermediate feature maps,
so its grouping network has no architectural dependency on the model
that produced those keypoints.
More importantly, this property is exercised experimentally rather than
inferred from the interface alone.
The detector-swap evaluation of Section~\ref{sec_res_e2e_swap} applies
the same frozen \texttt{meng-headline} checkpoint and frozen K-head
zero-shot to HigherHRNet,
YOLO11x-pose,
and RTMO-L outputs.
No grouping retraining is performed for either of the swapped detectors.

The autonomous PG-GAT pipeline remains close to each detector's native
person assignment,
the measured gaps are $0.0203$ for HigherHRNet,
$0.0210$ for YOLO11x-pose,
and $0.0220$ for RTMO-L.
The consistency of these values across three substantially different
detector families is the stronger result.
It shows not only that detector swapping is possible,
but that the accuracy cost remains approximately $0.02$ PGA across the
detectors tested.

The third clause is therefore supported in both respects.
The pipeline can estimate the person count from the detected evidence
without a ground-truth dependency,
and the same frozen grouping components can be used across multiple
detector families without retraining.
The autonomous component goes slightly beyond the original expectation -
predicted $K$ does not merely approximate the oracle-$K$ result under
real detector outputs,
but can provide a better clustering constraint when detection
incompleteness makes the annotation-level person count misaligned with
the evidence available to the grouping stage.